% Generated by code/rebuild/WriteResults.py from corrected-run outputs.
\section{Results}

To explore how a persistent change in investor decisions could destabilize
a market, we changed one response in selected traders' rules and found
that a Sell majority formed when enough traders received the new rule
(Fig.~\ref{fig:rule-quench}; Supplementary Tables~\ref{tab:campaigns}
and~\ref{tab:paired-switch}).
Rule~19286 changed the response to three neighbors in Buy and one in Hold
from Buy to Sell, preserving every other response
(Fig.~\ref{fig:rule-quench}A,B). Each treated market shared its baseline
history and first post-switch quote with a comparison market that retained
Rule~19280. With active technical traders and dealer depth $h=100$, none
of the 50 starting markets reached a majority at $f=0.10$ or 0.15, whereas
all reached one at $f=0.20$ and 0.25. No matched control reached a majority
over the same period. The median onset among the two larger fractions
changed from 1,070.0 to 311.5
steps (Supplementary Table~\ref{tab:paired-switch}). In the selected
pair, groups of sellers expanded after the rule change while the
comparison market retained mixed decisions (Fig.~\ref{fig:rule-quench}C).
The lasting rule change also differed from a temporary instruction:
setting the selected traders to Sell once at $f=0.15$, 0.20, or 0.25
produced no majority in either technical-trader condition (Supplementary Table~\ref{tab:revision-pulse}).
Among the eighteen single-response rule changes tested at $f=0.20$,
6 produced a majority in every tested starting market under both feedback conditions, while 12 produced none (Supplementary Table~\ref{tab:revision-rules}).
Within these tested interventions, majority formation differed with the
changed response, the fraction of traders receiving it, and whether the
intervention changed a rule or only the current decision.

To determine whether price feedback was required for widespread selling,
we repeated the rule switch while keeping technical traders in Hold and
found that a majority still formed (Fig.~\ref{fig:reservoir}).
With these traders held, majority probability was 0.34
at $f=0.175$, 0.70 at $f=0.185$, and
1.00 at $f=0.20$ (Fig.~\ref{fig:reservoir}A).
Every held market was simulated again with its full cash, share, and price
history. Across dealer depths, onset, final mean Sell fraction, and final
interior states were identical for the same starting market and target
fraction, although the associated financial paths differed. This agreed
with the held update, which had no input from financial variables. Neighbor responses
were therefore sufficient to produce a Sell majority. Allowing
technical traders to respond to price changed majority frequency near
the crossover between mostly mixed and mostly Sell-majority outcomes.
At $h=100$, the active condition gave 25 of 50
majorities at $f=0.175$ and 42 at $f=0.185$,
compared with 17 and
35 when held (Fig.~\ref{fig:reservoir}A,B).
Both the size and direction of the active-minus-held difference depended
on the tested depth and fraction. A price response could thus modify
how often the rule change produced a majority, but it was not necessary
for selling decisions to spread. These comparisons separated the local
decision mechanism from feedback through completed trades.

To determine how the spread of selling depended on the number of
interacting traders, we compared grid sizes with technical traders held
and found a sharper majority response on larger grids
(Fig.~\ref{fig:finite-size}). We refined the fractions near each
response while retaining the same starting markets. At the original
3,000-step horizon, the interpolated fraction at which half the markets
reached a majority, $f_{50}$, changed from 0.216 at $L=32$ to
0.166 at $L=128$. The width separating majority probabilities
of 0.25 and 0.75 changed from 0.023 [0.018, 0.027] to
0.008 [0.004, 0.010], with brackets giving 95\% intervals from
resampling whole starting markets (Fig.~\ref{fig:finite-size}A,C;
Supplementary Table~\ref{tab:finite-size-scaling}). Variation in the
final mean Sell fraction also changed: its variance multiplied by the
number of interior traders, the susceptibility, had sampled maxima of
21.2 and 241.2 at the
endpoint sizes (Fig.~\ref{fig:finite-size}B,D). Across all size and
feedback conditions, the largest absolute difference between logistic
and interpolated midpoints was 0.0035, indicating the sensitivity
of the estimate to the imposed curve shape
(Supplementary Tables~\ref{tab:finite-size-scaling} and~\ref{tab:revision-fits}). Observation time
remained consequential. At $L=64$, extending the held markets from
3,000 to 9,000 steps changed the majority midpoint from
0.178 to 0.167; alternative
Sell thresholds and size-dependent horizons were measured separately
(Supplementary Table~\ref{tab:revision-horizons}). The size response
therefore described a crossover within the observation window, while
the continued movement with time prevented its interpretation as a
fixed threshold for indefinitely long runs.


To examine how the spatial patterns developed beyond the initial
observation window, we extended five numerically selected markets and
found that longer observation changed some outcomes
(Fig.~\ref{fig:dose-panels}). Selection used the newly computed
3,000-step results and preceded inspection of images. Two selected
markets had not reached a majority at that endpoint, while three had.
At 10,000 steps, 1 of the initially nonmajority examples had
crossed one-half. The first-majority outcomes, in displayed order, were
$f=0.150$, starting market 2: not within 10,000 steps; $f=0.175$, starting market 15: after 4{,}073 steps; $f=0.175$, starting market 41: after 2{,}330 steps; $f=0.185$, starting market 5: after 1{,}836 steps; $f=0.200$, starting market 38: after 1{,}028 steps (Supplementary Table~\ref{tab:market-panel-horizons}).
The spatial histories connected these crossings to the growth and
merger of groups of sellers, while the price histories measured their
financial consequences (Fig.~\ref{fig:dose-panels}A--C). In the selected
$f=0.175$ market that had not crossed by 3,000 steps, maximum drawdown
changed from 10.17\% to
26.34\% at the longer endpoint
(Supplementary Table~\ref{tab:market-panel-horizons}).
We also extended all 50 starting markets at $f=0.150$, retaining
matched no-switch comparisons. The number of majorities among changed-rule
markets was 0 at 3,000 steps and
2 at 10,000 steps,
while 0 controls reached a majority by
10,000 steps. Median drawdown changed from
5.04\% to 5.87\%,
and the largest observed drawdown changed from
12.29\% to 25.30\%
(Supplementary Table~\ref{tab:market-panel-horizons}).
The longer trajectories therefore qualified outcomes defined by the
original endpoint; persistence of mixed decisions over that shorter
window did not establish persistence over longer times.

We varied dealer depth while retaining the same decision-rule intervention
and found that shallower positive reserves accompanied larger price
declines, while majority probability did not vary monotonically with
depth (Fig.~\ref{fig:reservoir}). At $f=0.20$, every starting market
reached a majority at every tested depth. With active technical traders,
median maximum drawdown was 99.99\% at $h=0.1$,
99.52\% at $h=1$, and 25.15\% at $h=100$
(Fig.~\ref{fig:reservoir}C). The corresponding median rejection
fractions were 0.000226, 0.000034, and
0.0000006, so small positive reserves supported large declines
while rejecting few orders backed by trader resources
(Fig.~\ref{fig:reservoir}D). Zero depth supplied a separate limit:
only opposing trader orders matched, net dealer trading was zero, and
price stayed fixed. No-switch markets at $h=0.1$ already had median
drawdown 77.04\%, compared with
0.75\% at $h=100$. The median paired
increase after the rule switch was 22.96
percentage points at $h=0.1$ and 24.39
at $h=100$ (Supplementary Table~\ref{tab:paired-switch}).
Positive reserve depth changed both available
funding and the price required to execute a net trade. In the additional
comparison at $h=40$ and $h=400$, median drawdowns were
48.62\% and
7.08\%, respectively
(Supplementary Table~\ref{tab:revision-impact}).
Near the majority crossover, depth also changed feedback through the
technical traders. At $f=0.175$, majority probabilities were
0.44 at $h=0.1$ and 0.50 at $h=100$.
Their difference changed from -0.06 with
adaptive band widths to -0.26 with frozen widths
(Supplementary Table~\ref{tab:revision-band}). The paired difference
between those contrasts was 0.20, with a 95\%
bootstrap interval [-0.08, 0.50], so its sign was
not resolved (Supplementary Table~\ref{tab:revision-band}). The selected comparison
from starting market 2 also reached a majority at all three
displayed depths while following distinct price paths
(Fig.~\ref{fig:dealer-panels}). The reserve mechanism thus separated
majority formation from the severity of the decline, but shallow positive
reserves did not prevent severe drawdowns.

To distinguish trader resources from dealer reserves, we changed cash
and share holdings independently at the rule switch and found that
financial resources changed interior decisions only through active
price feedback (Fig.~\ref{fig:balance-sheet}). At $f=0.18$, all
cash and share combinations with technical traders held produced the
same decision history within a starting market, giving majority
probability 0.64 across dealer depths. With technical traders
active, probabilities ranged from 0.54
to 0.80
(Fig.~\ref{fig:balance-sheet}A; Supplementary Table~\ref{tab:balance-fractions}).
Holdings entered the order and clearing calculations but did not enter
the interior lookup table. Their effect on subsequent decisions therefore
required the intervening price response. With unchanged cash and $h=100$,
share multipliers 0.1, 1, and 10 gave median drawdowns of
3.17\%, 20.18\%, and 77.55\%,
respectively (Fig.~\ref{fig:balance-sheet}B). More shares supplied
more potential sales, while the reserve curve determined the price
at which net selling could clear. Increasing the cash multiplier from 1
to 10 at $h=1$ with unchanged shares changed median drawdown from 96.46\%
to 97.10\%, and majority probability from
0.58 to 0.56
(Fig.~\ref{fig:balance-sheet}A,B). Rejection recorded the fraction
of orders backed by trader resources that still failed to execute,
and its relationship to drawdown differed across the resource
combinations (Fig.~\ref{fig:balance-sheet}C,D). The financial
consequences of a rule change therefore depended on both sides of
traders' books and the dealer's reserves. A majority decision to sell
alone did not specify how many shares could change hands or how far
the marginal quote would fall.
